\section{Correctness}
\label{sec:correctness}
The arguments concern the core protocol, including the prefix-wide overlap guard, late notarization, and Rule~3 promotion, and they are relative to the services of Section~\ref{sec:services}: safety uses C1 and A1--A3, termination additionally uses C2 and A4. They do not include the optional fresh-work extension of the supplement. Cryptographic guarantees hold except with negligible probability.

\subsection{Safety}
\begin{lemma}[Account exclusion]
\label{lem:exclusion}
In one account instance, at most one block has an $\NC$, whether it counts first votes, late notarizations, or both. An $\FC$ or $\FF$ excludes a conflicting $\NC$. If an $\FF$ is formable, every rival has at most $f+p<r$ first-vote supporters.
\end{lemma}
\begin{proof}
Two $q$-quorums intersect in at least $2q-N=f+1$ identities, one of them correct, which supports one block only in the instance, by a first vote or a late notarization (Section~\ref{sec:voting}). An $\FC$ has a correct final voter who first obtained an $\NC$; an $\FF$ contains $N-p\geq q$ first votes. For the last claim, with actual Byzantine count $f'\leq f$, an $\FF$ has at least $N-p-f'$ correct first voters, so at most $p$ correct identities remain for rivals; adding every Byzantine identity gives at most $p+f'\leq p+f$.
\end{proof}

\begin{lemma}[Late notarization adds no old finality]
\label{lem:late}
Every epoch-$e$ $\FC$ or $\FF$ that can be formed consists of correct votes released before their signers froze, together with Byzantine votes.
\end{lemma}
\begin{proof}
A correct validator releases no epoch-$e$ first or final vote after leaving epoch $e$ (Section~\ref{sec:transitions}), and neither certificate counts a late notarization: an $\FF$ counts first votes, an $\FC$ final votes.
\end{proof}

\begin{lemma}[Frozen-report exposure]
\label{lem:exposure}
A correct reporter's released final vote implies a frozen $\tagN$ or $\tagF$ entry for its selected block. Its first vote implies a frozen ledger entry for that block or a membership in $G_i$. These statements persist after report creation.
\end{lemma}
\begin{proof}
Final voting follows durable eligible notarization, and account exclusion prevents incompatible same-domain finality from removing the selected path. $G_i=V_i^{(e)}\setminus\keys(T_i)$ partitions the reporter's old-domain voted hashes. Stopping precedes freezing; roots and tags are then immutable, and C1 ensures that what a correct validator reconstructs is exactly what was frozen.
\end{proof}

\begin{theorem}[Closure preserves latent and late finality]
\label{thm:closure}
Fix a valid closed predecessor. If an eligible old-epoch account final certificate can be formed, every eligible candidate (every $v$ with $\Elig_e(v)$) preserves its selected branch and ancestors, as compatible finality or as the locks that certificate requires. The certificate may be assembled after report freeze or after the checkpoint is decided.
\end{theorem}
\begin{proof}
Let $R$ be the selected reporters. For a possible normal certificate with final-voter set $A$, $|(R\cap A)\setminus\Byz|\ge q-2f=p+1$; each such correct reporter froze an $\tagN$ or $\tagF$ entry, where $\tagF$ exposes finality and $\tagN$ triggers Rule~1. For a possible fast certificate, $|(R\cap A)\setminus\Byz|\ge(N-p)-2f=r$; each of these reporters exposes the block through $T_i$ or through $G_i$ with its verified first vote, and in the latter case Rule~2 retains it. Lemma~\ref{lem:exclusion} prevents valid discharge or a threshold rival from eliminating genuine fast finality. Usability and the committed manifest make these memberships available to $\Build$. Later unselected or Byzantine signatures do not remove the required correct intersections.
\end{proof}

\begin{lemma}[Checkpoint-promotion safety]
\label{lem:promotion}
A Rule~3 promoted block cannot conflict with an eligible closing-epoch $\NC$, $\FC$, or $\FF$, so a decision that finalizes the promoted prefix is compatible with old-domain finality and with core overlap finality.
\end{lemma}
\begin{proof}
Every promoted block has a verified closing-epoch $\NC$; Lemma~\ref{lem:exclusion} excludes a conflicting old $\NC$ and hence conflicting old $\FC/\FF$. Core overlap finality for a conflicting successor block would require a conflicting closing-epoch $\NC$ at the first divergent position, also impossible. Promotion stops at retained forks and never promotes recovery-only evidence.
\end{proof}

\begin{lemma}[Acyclic final dependencies]
\label{lem:dag}
For any finite set of valid final blocks closed under account ancestors and required finalized sends, parent and send-to-receive edges form a directed acyclic graph.
\end{lemma}
\begin{proof}
Order blocks by the earliest event at which valid final evidence justifies them. Along a parent edge this event cannot become earlier at the child; along a send-to-receive edge it is strictly later at the receive, because admission required the send's finality proof and Rule~3 cannot supply it for the checkpoint under construction. A cycle cannot consist only of parent edges, since slots increase strictly, so it contains a send-to-receive edge and a strict increase around a cycle, a contradiction.
\end{proof}

\begin{lemma}[Deterministic construction]
\label{lem:construction}
Assume the predecessor satisfies the closure invariant. For fixed usable selected reports and a valid immutable manifest, $\Build$ returns one canonical candidate state at every correct validator. It preserves inherited finality, carries all mandatory surviving locks, and applies only explicit account finality plus Rule~3 promotion. All-correct usable inputs yield a finite valid construction.
\end{lemma}
\begin{proof}
By C1 every correct validator reconstructs the same $(T_i,G_i)$ for each selected header, and the manifest fixes witnesses and proof choices; hence explicit finality and $R_Q$ are the same everywhere. Lemma~\ref{lem:exclusion}, predecessor validity, attachment, and Theorem~\ref{thm:closure} make the retained finality and locks compatible. Rule~3's tests fix $P_Q$, safe by Lemma~\ref{lem:promotion}; Lemma~\ref{lem:dag} and the canonical tie-break fix replay. Malformed Byzantine inputs are excluded by usability; all-correct usable inputs yield a finite state.
\end{proof}

\begin{theorem}[Transfer safety across epochs]
\label{thm:safety}
Assuming C1 and A1--A3, finalized histories at correct validators are compatible and irreversible across all core-protocol epochs; finalized balances are nonnegative, and each send is received at most once.
\end{theorem}
\begin{proof}
Induct from genesis. Within an epoch, Lemma~\ref{lem:exclusion} excludes incompatible account finality. At a boundary, A2 ensures that the decided $v_e$ is eligible, so by Theorem~\ref{thm:closure} it preserves every formable old final branch and by Lemma~\ref{lem:promotion} its promotion is compatible; A1 and A3 ensure that every correct validator installs the state of the same $v_e$, which by Lemma~\ref{lem:construction} is the same state. For overlap, the first divergent unresolved position on an early successor-final prefix has a closing-epoch $\NC$ for the successor block by the prefix-wide guard, possibly assembled from late notarizations, which by Lemma~\ref{lem:late} add no old finality that Theorem~\ref{thm:closure} would have to preserve; a conflicting closing-epoch certificate or a conflicting Rule~3 promotion would require a conflicting $\NC$, contradicting Lemma~\ref{lem:exclusion}. Matching-origin discharge and inherited locks remain enforced after installation, and the base/overlay union with applied-operation identifiers prevents rollback and duplicate application. Replay checks balances and finalized-send dependencies.
\end{proof}

The proof uses nothing about how the consensus service reaches its decision: an eligible value is safe whichever one is chosen, because Theorem~\ref{thm:closure} is universal over eligible candidates. This is what lets RAI leave the service abstract.

\subsection{Conditional termination}
\begin{lemma}[Correct reports become usable]
\label{lem:reports-live}
Under the data service, finite relevant admitted work, fair processing, and C2, every correct report eventually becomes usable at every correct validator.
\end{lemma}
\begin{proof}
A correct reporter stops old first and final voting, so its old vote set and roots are finite, and its report is well formed. By C2 every correct validator eventually reconstructs $(T_i,G_i)$. The data service delivers the blocks, ancestry, signed votes, and witnesses behind its claims, and fair processing completes the finite validation. A withholding Byzantine reporter can make only its own report unusable; $N-f$ correct reports remain.
\end{proof}

\begin{theorem}[Handoff termination]
\label{thm:handoff-live}
Assume the data service, fair processing, C2, A4, and recurring prepared windows with $N-f$ responsive correct members of $K_e$ long enough for reconstruction, candidate preparation, and the consensus service's rounds. Every enabled handoff is eventually decided and, with eventual retained-data delivery, installed by correct successors.
\end{theorem}
\begin{proof}
By Lemma~\ref{lem:reports-live}, $N-f$ correct reports become usable at every correct member, which by Lemma~\ref{lem:construction} can then construct an eligible candidate; the canonical selection order makes correct members that hold the same usable reports construct the same one. A4 then decides some eligible value in a sufficient window; A3 lets successors verify it, and retained authorizing data let them recompute and install the state, which enables the next deferred boundary.
\end{proof}

\begin{theorem}[Uncontested account continuation]
\label{thm:account-live}
Assume handoff termination, available final-send dependencies, and recurring common-open ready intervals. An owner following the retry-or-fresh-child policy of Section~\ref{sec:account} eventually finalizes each target of a finite valid uncontested continuation, provided each required new position has such an interval, and the finalized targets eventually enter a checkpoint.
\end{theorem}
\begin{proof}
In a ready interval, $N-f\ge q$ correct first votes create an $\NC$ and the same validators can final-vote. An attempt interrupted by a boundary is preserved by Theorem~\ref{thm:closure}; an omitted block can be retried, a uniquely retained notarized target may finalize at the handoff, and a remaining retained fork can be selected by a fresh uncontested child. Repeating this over the finite continuation and using eventual evidence delivery gives finalization and checkpoint inclusion.
\end{proof}

With ancestry prepared, normal finality takes three delivery stages after broadcast: block, first votes, final votes; an available $N-p$ first-vote fast quorum takes two. A sufficient account interval budget is
\begin{equation}
\frac{D}{1+\rho}>\kappa_A+R_A+3\delta,\label{eq:account-window}
\end{equation}
where $\kappa_A$ bounds opening skew and $R_A$ bounds readiness after the last opening, including predecessor closure. Deferred boundaries only lengthen an open interval. The corresponding window for the consensus service depends on the instantiation; Section~\ref{sec:eval} gives the condition for ours. These are conditional communication bounds, not measured latency.
